\chapter{动机}

\section{动机领域}

动机（Motivation）维护主体面向未来状态的长期认知。基本对象包括愿望、目标、关切和优先关系。

\begin{description}
\item[愿望（Desire）] 主体希望某种状态出现或持续的倾向。
\item[目标（Goal）] 可被识别、跟踪并与结果对应的期望未来状态。
\item[关切（Concern）] 需要持续关注、保护或避免恶化的问题和状态。
\item[优先关系（Priority）] 多个动机对象在特定时间或条件下的重要性关系。
\end{description}

动机对象具有来源、修订和生命周期。具体的满足、放弃、冲突、暂停和恢复状态由下游动机规格细化；这些状态不得改变认知目标与行为承诺的系统边界。

\section{目标与承诺}

Nous Wave 的目标表示主体希望实现的未来状态。Heptalogos 的承诺、义务、计划和任务表示已经进入行为或规范体系的约束性状态。

目标可以被 Heptalogos 采纳、延迟、重述或拒绝。采纳形成行为承诺时，行为侧保留来源目标及其修订引用。执行结果、完成、失败、取消或拒绝通过类型明确的结果回执返回 Nous Wave，并可触发动机或其他认知领域的修订提案。

记忆、自我或社会认知中的一句自然语言陈述不能直接形成外部承诺。任何行为承诺都经过 Heptalogos 的行为与规范审查。

\section{动机与运行时}

持久动机保存长期愿望、目标和关切；认知运行时决定当前哪些动机进入前台并影响注意、检索和即时评价。运行时优先级可以随情境变化，不要求每次前台调整都形成新的持久动机修订。
